\section{Données et méthodologie}
\label{sec:methodo}

Cette section présente les données et la méthodologie de la recherche. La
pauvreté multidimensionnelle des enfants est mesurée par l'approche N-MODA de
l'UNICEF, qui sert de variable de résultat.
L'impact causal des transferts est estimé sur les deux vagues de l'EHCVM par la
combinaison du score de propension (PSM) et de la double différence (DD).

% ─────────────────────────────────────────────────────────────────────────────
\subsection{Sources de données}
% ─────────────────────────────────────────────────────────────────────────────

\subsubsection{Présentation des enquêtes}

Les Enquêtes Harmonisées sur les Conditions de Vie des Ménages (EHCVM) sont
réalisées dans le cadre du programme statistique de l'UEMOA
\parencite{ansd2019, ansd2022}. Les deux vagues mobilisées dans ce travail
reposent sur un plan de sondage aléatoire stratifié à deux degrés :

\begin{itemize}
  \item Strates : les 14 régions administratives du Sénégal sont chacune
  subdivisées suivant le milieu de résidence (urbain, rural), soit 28 strates.
  Le tirage est indépendant dans chaque strate, ce qui garantit la
  représentativité au niveau région et milieu de résidence;
  \item Unités primaires : des districts
  de recensement (DR) issus du RGPHAE 2013 constituent les grappes.
  598 DR ont été sélectionnés (330 urbains, 268 ruraux);
  \item Unités secondaires : 12 ménages
  sont enquêtés dans chaque DR.
\end{itemize}

\subsubsection{Construction de l'échantillon}

Notre population cible est l'ensemble des enfants âgés de 0 à 17 ans résidant
dans les ménages enquêtés dans les deux vagues. L'EHCVM I couvre 7\,156 ménages et l'EHCVM II 7\,120 ménages selon le même plan de sondage avec un
sous-échantillon panel identifié dans les données par un indicateur de statut
de suivi \parencite{ansd2019, ansd2022}.

L'échantillon analytique repose sur le panel vrai constitué des
6\,127 ménages panel, c'est-à-dire les
ménages enquêtés lors des deux vagues et retrouvés dans leur logement d'origine
(5\,643 ménages dans le même logement ; 290 ménages
retrouvés dans un nouveau logement). Chaque ménage est ainsi observé en $t = 0$
(2018-2019) et $t = 1$ (2021-2022), ce qui autorise une estimation DD en
différences intra-ménage sans recourir à des hypothèses d'agrégation au niveau
des grappes \parencite{deaton1997}.

Le tableau~\ref{tab:panelhh} présente la répartition des ménages enquêtés
en 2021-2022 selon leur statut de suivi et le milieu de résidence.

\begin{table}[H]
  \centering
  \caption{Répartition des ménages EHCVM II selon le statut de suivi}
  \label{tab:panelhh}
  \small
  \zebre
  \begin{tabular}{|l|r|r|r|r|r|}
    \hline
\entete & \multicolumn{2}{|c|}{Nouveau ménage} & \multicolumn{2}{|c|}{Ménage panel} & Total \\ \hline
    
    & $n$ & \% & $n$ & \% & $n$ \\
    \hline
    Urbain & 740 & 18,9 & 3\,182 & 81,1 & 3\,922 \\ \hline
    Rural  & 253 &  7,9 & 2\,945 & 92,1 & 3\,198 \\
    \hline
    Total & 993 & 13,9 & 6\,127 & 86,1 & 7\,120 \\
    \hline
    \multicolumn{6}{|c|}{\textit{Dont ménages panel selon condition de logement}} \\ \hline
    \quad Même logement qu'en 2018 & & & 5\,643 & 92,1 & \\ \hline
    \quad Logement différent       & & &   290 &  4,7 & \\ \hline
    \quad Non renseigné            & & &   194 &  3,2 & \\
    \hline
  \end{tabular}
  \begin{tablenotes}
    \small
    \item \textit{Source :} EHCVM II, calcul de l'auteur
  \end{tablenotes}
\end{table}

Après nettoyage (exclusion des observations avec plus de deux indicateurs de
privation manquants et des ménages sans information sur les transferts),
l'échantillon retenu comprend les enfants âgés de 0 à 17 ans pour lesquels
l'intégralité des indicateurs N-MODA ainsi que la variable de traitement
sont disponibles dans les deux vagues. L'échantillon complet incluant les
993 nouveaux ménages de 2021 est utilisé dans
les analyses de robustesse. La limite principale de cette construction est
l'attrition différentielle : si les ménages perdus entre les deux vagues
diffèrent systématiquement des ménages retenus, les estimations PSM-DD
peuvent être biaisées, ce point est discuté dans la section~\ref{sec:discussion}.

\subsection{Construction de la variable de traitement}

La variable de traitement $D_i \in \{0, 1\}$ est un indicateur stable,
fixe au niveau du ménage sur les deux périodes, qui vaut 1 si le ménage de
l'enfant $i$ a reçu des transferts d'un membre émigré à l'étranger au cours
des 12 mois précédant l'enquête dans chacune des deux vagues (2018
et 2021). Cette définition garantit la stabilité du statut de
traitement entre les deux vagues, condition indispensable à l'estimateur de
double différence, qui requiert un indicateur de groupe fixe dans le temps, et est conforme aux conventions de la littérature empirique sur les
remittances en Afrique \parencite{azam2016, combes2019}.

La construction de cette variable repose sur la section 13 de transferts des
deux vagues EHCVM.

Un ménage est classé comme bénéficiaire stable ($D_i=1$) s'il satisfait les
deux conditions suivantes :


(i) il a reçu en 2018 au moins un transfert dont
l'expéditeur est localisé hors du Sénégal ;

(ii) il a également reçu un tel transfert en 2021.

Le groupe témoin ($D_i = 0$) est composé de ménages n'ayant reçu aucun
transfert international ni en 2018 ni en 2021 (never treated).

Les ménages dont le statut change entre les deux vagues, entrants tardifs
($D=0$ en 2018 puis $D=1$ en 2021) et sortants ($D=1$ puis $D=0$), sont
exclus de l'analyse causale, leur inclusion biaisant l'estimateur de la
double différence \parencite{callaway2021}. Sur les 6\,127 ménages du panel
vrai, cette classification donne 681 ménages traités stables, 3\,981 ménages
jamais traités et 1\,465 switchers exclus (816 entrants tardifs et
649 sortants). Les transferts internes sont exclus dans les deux vagues car ils
relèvent d'une logique de redistribution domestique distincte des envois
de la diaspora internationale \parencite{ndiaye2013}.

Une définition alternative du traitement aurait pu retenir les entrants
tardifs comme groupe traité (comparés aux jamais bénéficiaires), un design
de double différence plus classique où le traitement démarre entre les deux
vagues. Ce choix a été écarté au profit du traitement stable pour trois
raisons. D'abord, le passage au statut de bénéficiaire coïncide le plus
souvent avec un départ en migration récent d'un membre du ménage, lui-même
motivé par un choc économique antérieur~; l'évolution observée capterait
alors en partie l'effet de ce choc initial plutôt que celui du transfert.
Ensuite, ce départ modifie mécaniquement la composition du ménage (taille,
main-d'\oe uvre disponible, rapport de dépendance) au moment même où l'issue
est mesurée, un facteur de confusion propre aux entrants et largement absent
chez les bénéficiaires stables, dont la migration est antérieure à la période
de base et dont le ménage s'est déjà réorganisé en 2018. Enfin, la date exacte
d'entrée dans le statut de bénéficiaire n'est pas observée à l'intérieur de la
fenêtre 2019-2021, si bien que l'intensité et la durée du traitement varient
fortement au sein du groupe des entrants, contrairement au traitement stable,
observé aux deux bornes de la période. Cette définition alternative du
traitement fait néanmoins l'objet d'une estimation de robustesse, rapportée
à titre de comparaison en annexe~\ref{ann:A}.

% ─────────────────────────────────────────────────────────────────────────────
\subsection{Mesure de la pauvreté multidimensionnelle des enfants : approche MODA}
% ─────────────────────────────────────────────────────────────────────────────

L'approche MODA, développée par \textcite{deneubourg2012} au sein du bureau
de recherche de l'UNICEF et appliquée en afrique subsaharienne par
\textcite{demilliano2014}, part du constat que les indices synthétiques
classiques de pauvreté multidimensionnelle présentent deux limites lorsqu'ils
sont appliqués aux enfants. D'abord, ils utilisent les mêmes indicateurs et
seuils quel que soit l'âge, alors que les besoins fondamentaux varient
considérablement entre un nourrisson et un adolescent. Ensuite, ils héritent
de la tradition centrée sur le ménage et ne reflètent pas directement les
privations de l'enfant individuel \parencite{roelen2008}, les poids et le seuil
inter-dimensionnel étant en outre fixés de manière normative plutôt que dérivés
des configurations de privations effectives.

MODA répond à ces limites selon trois principes. Premièrement, l'enfant
individuel, et non le ménage, est l'unité d'analyse. Deuxièmement, les
indicateurs sont désagrégés par groupe d'âge pour refléter les droits
spécifiques à chaque stade de développement, tels qu'ils sont consacrés par les
instruments internationaux relatifs aux droits de l'enfant, au premier rang
desquels la Convention internationale des droits de l'enfant \parencite{onu1989}
et la Charte africaine des droits et du bien-être de l'enfant
\parencite{oua1990}~: chaque
dimension du N-MODA correspond à un droit fondamental garanti par ces
textes~: selon l'article 24 de la CIDE et l'article 14 de la Charte
africaine, la santé et la nutrition relèvent du droit à la survie et au
meilleur état de santé possible~; selon les articles 28 et 29 de la CIDE et
l'article 11 de la Charte africaine, l'éducation relève du droit à
l'instruction~; selon l'article 27 de la CIDE, l'eau, l'assainissement et le
logement relèvent du droit à un niveau de vie suffisant~; et selon les
articles 7, 9 et 32 de la CIDE ainsi que les articles 6 et 15 de la Charte
africaine, la protection de l'enfant relève du droit d'être préservé de
toute forme d'exploitation, de séparation abusive d'avec ses parents et à
l'enregistrement à la naissance \parencite{onu1989,oua1990}, ce qui
ancre la mesure dans une définition normative du bien-être de l'enfant
plutôt que dans le
seul revenu du ménage. Troisièmement, l'analyse
porte sur les privations superposées (overlapping deprivations) :
quelles privations se combinent chez les enfants les plus vulnérables, et
avec quelle intensité.

L'analyse par privations superposées fournit ainsi une lecture fine de la
structure des privations par âge et de leurs combinaisons. Appliquant cette
méthodologie à trente pays d'Afrique subsaharienne, \textcite{demilliano2014}
montrent que les configurations de privations varient considérablement selon le
groupe d'âge et le pays, ce qui justifie empiriquement l'approche désagrégée~:
en Afrique de l'Ouest, les privations les plus fréquentes chez les jeunes
enfants (0-4 ans) concernent l'accès à l'eau potable et l'assainissement,
tandis que pour les enfants d'âge scolaire (5-14 ans), la non-scolarisation et
le retard scolaire occupent une place centrale. Cette dépendance des dimensions
au groupe d'âge est précisément ce que mobilise notre analyse de robustesse par
panel d'enfants (section~\ref{subsec:robustesse})~: un enfant suivi entre
2018 et 2021 vieillit d'environ trois ans et peut ainsi changer de groupe
d'âge, de sorte que l'ensemble des dimensions qui lui sont applicables se
modifie d'une vague à l'autre.


Suivant \textcite{deneubourg2012} et l'adaptation nationale N-MODA développée
pour le Sénégal lors d'un atelier participatif rassemblant l'ANSD et les
ministères sectoriels \parencite{ansd2024}, nous distinguons trois groupes
d'âge avec des indicateurs adaptés aux droits spécifiques de chaque phase de
l'enfance. La matrice complète des dimensions, indicateurs, seuils de privation
et de leur applicabilité par groupe d'âge est détaillée en
annexe~\ref{ann:C} (tableau~\ref{tab:matrice_nmoda}).
Les sept dimensions retenues sont l'assainissement, l'eau, le logement, la
nutrition, la santé, la protection de l'enfant et l'éducation, chacune reposant
sur des indicateurs comparables entre les deux vagues de l'EHCVM.

Pour chaque enfant $i$ et chaque dimension $j \in \{1, \ldots, 7\}$, la variable
de privation est $d_{ij} = \mathbf{1}[\text{privé dans au moins un indicateur
de la dimension } j]$. Le \textit{score de
privation} de l'enfant $i$ est le nombre de dimensions dans lesquelles il est privé :
\begin{equation}
  s_i = \sum_{j=1}^{7} d_{ij}.
  \label{eq:score_moda}
\end{equation}

Le taux de privation par dimension mesure la prévalence de chaque
privation :
\begin{equation}
  H_j = \frac{1}{n} \sum_{i=1}^{n} d_{ij}.
  \label{eq:hj}
\end{equation}

Le taux de pauvreté multidimensionnelle $H$ (incidence) est la part
d'enfants privés dans au moins $k$ dimensions simultanément :
\begin{equation}
  H = \frac{1}{n} \sum_{i=1}^{n} \mathbf{1}\left[s_i \geq k\right].
  \label{eq:pk}
\end{equation}
Conformément à la méthodologie N-MODA Sénégal \parencite{ansd2024}, nous
retenons qu'un enfant est considéré multidimensionnellement pauvre s'il
est privé dans au moins 4 des 7 dimensions simultanément ($k = 4$).

L'intensité moyenne des privations parmi les enfants pauvres est :
\begin{equation}
  A = \frac{1}{q} \sum_{i:\,s_i \geq k} \frac{s_i}{7},
  \label{eq:intensite_moda}
\end{equation}
où $q$ est le nombre d'enfants pauvres. L'indice synthétique $M_0 = H \times A$
mesure la pauvreté multidimensionnelle ajustée pour l'intensité.

L'approche N-MODA fournit ainsi la variable d'impact : l'indicateur
binaire identifiant le statut de pauvreté multidimensionnelle de l'enfant (privé dans au moins $k=4$ dimensions sur 7). Les
taux de privation par dimension $H_j$ et l'indice ajusté $M_0 = H \times A$
servent de mesures secondaires, décomposables par milieu, région et groupe
d'âge.

% ─────────────────────────────────────────────────────────────────────────────
\subsection{Stratégie d'identification causale}
% ─────────────────────────────────────────────────────────────────────────────

\subsubsection{Cadre des résultats potentiels}

Pour chaque enfant $i$, on définit deux résultats potentiels : $Y_i(1)$ la
mesure de pauvreté N-MODA si le ménage reçoit des transferts, et
$Y_i(0)$ si le ménage n'en reçoit pas. On observe seulement $Y_i = D_i \cdot
Y_i(1) + (1 - D_i) \cdot Y_i(0)$. Le paramètre d'intérêt est l'effet moyen
du traitement sur les traités (ATT) :
\begin{equation}
  \tau_{ATT} = \E\left[Y_i(1) - Y_i(0) \mid D_i = 1\right].
  \label{eq:att}
\end{equation}
L'estimateur naïf est biaisé car les ménages bénéficiaires et non-bénéficiaires
diffèrent systématiquement sur des caractéristiques observables et inobservables
(biais de sélection).

\subsubsection{Appariement par score de propension (PSM)}

L'appariement par score de propension, introduit par \textcite{rosenbaum1983},
corrige le biais de sélection lié aux caractéristiques observables dans les
études non expérimentales. Plutôt qu'une comparaison directe entre ménages
bénéficiaires et non-bénéficiaires, potentiellement très différents dès le
départ, il construit pour chaque ménage bénéficiaire un contrefactuel crédible
à partir de ménages non-bénéficiaires aux caractéristiques observables
similaires. Le PSM repose sur deux hypothèses. La première est l'indépendance
conditionnelle (CIA) : conditionnellement aux covariables observées $X_i$,
les résultats potentiels sont indépendants du traitement,
\begin{equation}
  \left(Y_i(0),\, Y_i(1)\right) \perp D_i \mid X_i.
  \label{eq:cia}
\end{equation}
La seconde est le support commun : $0 < P(D_i = 1 \mid X_i) < 1$,
garantissant que l'appariement est réalisable pour chaque traité.
\textcite{rosenbaum1983} ont montré que si la CIA est satisfaite pour $X_i$,
elle l'est aussi pour le score de propension $p(X_i) = P(D_i = 1 \mid X_i)$,
estimé par un modèle probit :
\begin{equation}
  p(X_i) = \Phi\left(X_i'\hat{\beta}\right).
  \label{eq:probit}
\end{equation}

Les covariables $X_i$ du modèle probit sont choisies pour satisfaire
l'hypothèse d'indépendance conditionnelle (CIA) tout en évitant d'inclure
des variables potentiellement affectées par le traitement~: chacune est
mesurée à la période de base (2018) et se rattache à l'un des cadres
théoriques exposés en revue de littérature (section~\ref{sec:revue}), qui
gouvernent conjointement la décision de migrer, d'envoyer des transferts et,
partant, le statut de bénéficiaire du ménage.

Les caractéristiques du chef de ménage relèvent d'abord de la théorie
néoclassique du capital humain \parencite{sjaastad1962, harris1970}~: l'âge et
le niveau d'éducation du chef influencent la capacité du ménage à financer et
organiser la migration d'un membre, tandis que le sexe du chef capte la
configuration migratoire typique où l'époux émigré laisse la gestion du
ménage à son épouse, ce qui explique la sur-représentation des chefs féminins
parmi les ménages bénéficiaires. Le statut matrimonial complète ce profil, la
structure familiale conditionnant à la fois la probabilité de migration d'un
membre et le motif altruiste des transferts vers le conjoint resté au pays
\parencite{lucas1985}.

Les caractéristiques de composition du ménage renvoient à la nouvelle
économie des migrations de travail \parencite{stark1985}~: la taille du
ménage détermine la main-d'\oe uvre familiale disponible pour libérer un
migrant sans compromettre les activités domestiques, tandis que le
bien-être initial (dépense par tête, transformée en logarithme) capte à la
fois la contrainte de liquidité qui conditionne le financement des coûts de
migration \parencite{sjaastad1962} et le motif de diversification du risque
qui incite les ménages les plus vulnérables à recourir à la migration comme
assurance de fait \parencite{stark1985, gubert2002}. Cette variable est en
outre indispensable pour satisfaire la CIA, le niveau de vie initial
déterminant conjointement le statut de traitement et la trajectoire ultérieure
de pauvreté des enfants.

Enfin, le milieu de résidence et la région renvoient à la théorie des réseaux
migratoires \parencite{massey1993, bourdieu1986}~: l'accès aux réseaux et aux
canaux de transfert varie fortement selon la localisation, certaines régions
sénégalaises (vallée du fleuve, Casamance) concentrant une tradition
migratoire ancienne qui abaisse les coûts et les risques de départ pour leurs
résidents, indépendamment de leurs caractéristiques individuelles. Cette
spécification rejoint celle retenue dans la littérature empirique sur les
déterminants des transferts en Afrique \parencite{adams2005, gubert2002,
azam2016, combes2019}.

\begin{table}[H]
  \centering
  \caption{Covariables du modèle de propension}
  \label{tab:covariables}
  \small
  \zebre
  \begin{tabular}{|l|l|}
    \hline
\entete Groupe & Variable \\
    \hline
    Chef de ménage & Âge du chef de ménage \\ \hline
                   & Sexe (féminin = 1)    \\ \hline
                   & Niveau d'éducation    \\ \hline
                   & Statut matrimonial    \\ \hline
    Composition    & Taille du ménage      \\ \hline
                   & Bien-être initial (dépense par tête) \\ \hline
    Localisation   & Milieu (urbain = 1)   \\ \hline
                   & Région (14 régions)   \\ \hline
  \end{tabular}
  \begin{minipage}{\linewidth}
    \small\vspace{2pt}
    \textit{Note :} le bien-être initial (dépense par tête) est transformé en
    logarithme pour réduire l'asymétrie de la distribution~; la région est
    traitée comme un facteur à 14 modalités. La spécification retient la taille
    du ménage, le logarithme de la dépense par tête, le milieu, la région, le
    sexe, l'âge, l'éducation et le statut matrimonial du chef de ménage.
  \end{minipage}
\end{table}

L'appariement est réalisé au niveau du ménage : toutes les covariables
du score de propension étant des caractéristiques du ménage,
tous les enfants d'un même ménage partagent le même score. Apparier au niveau
de l'enfant induirait des ex-aequo massifs et une pondération arbitraire par
le nombre d'enfants~; l'appariement au niveau ménage évite ces artefacts, les
poids d'appariement étant ensuite appliqués à l'ensemble des enfants du ménage
dans l'estimation. Ce poids n'est pas divisé par le nombre d'enfants du
ménage~: chaque enfant hérite du poids complet de son ménage, de sorte qu'un
ménage comptant davantage d'enfants pèse proportionnellement plus dans
l'estimation. Ce choix est cohérent avec l'unité d'analyse retenue,
l'enfant, l'objectif étant de mesurer l'effet moyen sur les enfants et non
sur les ménages.

Trois algorithmes d'appariement sont mis en œuvre et comparés \parencite{caliendo2008,
heckman1997}. L'appariement aux $k$ plus proches voisins \parencite{rosenbaum1983}
($k = 4$, avec remplacement) associe à chaque traité les 4 témoins les plus proches
en termes de score :
\begin{equation}
  \hat{\tau}_{k\text{-NN}} = \frac{1}{n_T} \sum_{i \in \mathcal{T}}
  \left[Y_i - \frac{1}{k}\sum_{j \in \mathcal{N}(i)} Y_j\right].
  \label{eq:knn}
\end{equation}
L'appariement par noyau d'Epanechnikov \parencite{heckman1997} utilise une moyenne
pondérée de l'ensemble des témoins, les poids décroissant avec la distance en score :
\begin{equation}
  \hat{\tau}_{\text{kernel}} = \frac{1}{n_T} \sum_{i \in \mathcal{T}}
  \left[Y_i - \frac{\sum_{j \in \mathcal{C}} K\!\left(\frac{p_i-p_j}{h}
  \right) Y_j}{\sum_{j \in \mathcal{C}} K\!\left(\frac{p_i-p_j}{h}\right)}
  \right],
  \label{eq:kernel}
\end{equation}
où $h = 0{,}06$ est la largeur de bande. L'appariement
par rayon (caliper, $\epsilon = 0{,}05$, sans remise) \parencite{caliendo2008}
exclut les appariements de faible qualité en imposant une distance maximale
sur le score de propension.

La qualité de l'appariement est évaluée par le biais standardisé
\parencite{rosenbaum1983},
\begin{equation}
  SB_j = 100 \cdot \frac{\bar{X}_{1j} - \bar{X}_{0j}}
  {\sqrt{0{,}5\left(V_{1j} + V_{0j}\right)}},
  \label{eq:sb}
\end{equation}
un $|SB| < 10$ traduit un équilibre satisfaisant \parencite{caliendo2008},
ainsi que par le test $t$ d'égalité des moyennes et la statistique pseudo-$R^2$.

\subsubsection{Double différence (DD)}

La DD exploite la dimension temporelle des deux vagues EHCVM pour contrôler
les effets fixes inobservables stables dans le temps :
\begin{equation}
  \hat{\tau}_{DD} = \left(\bar{Y}^1_{t=2} - \bar{Y}^1_{t=1}\right)
  - \left(\bar{Y}^0_{t=2} - \bar{Y}^0_{t=1}\right),
  \label{eq:dd}
\end{equation}
où les exposants 1 et 0 désignent respectivement le groupe traité ($D_i=1$)
et le groupe témoin ($D_i=0$), sous l'hypothèse de tendances parallèles :
\begin{equation}
  \E\!\left[Y_i(0)_{t=2} - Y_i(0)_{t=1} \mid D_i = 1\right]
  = \E\!\left[Y_i(0)_{t=2} - Y_i(0)_{t=1} \mid D_i = 0\right].
  \label{eq:parallel}
\end{equation}
Parce que $D_i$ est figé sur la valeur de 2018, les ménages dont le statut
de réception change entre les deux vagues sont exclus de l'échantillon
d'estimation (restriction never-treated pour le groupe témoin).
Cette exclusion garantit que le terme $\bar{Y}^0_{t=2} - \bar{Y}^0_{t=1}$
décrit bien la tendance contrefactuelle des non-bénéficiaires, non contaminée
par une réception tardive de transferts \parencite{callaway2021}.

\subsubsection{Estimateur PSM-DD}

La combinaison PSM-DD proposée par \textcite{heckman1997} et
\textcite{heckman1998} neutralise simultanément le biais lié aux différences
observables (PSM) et celui lié aux effets fixes inobservables constants dans
le temps (DD) :
\begin{equation}
  \hat{\tau}_{PSM\text{-}DD} = \frac{1}{n_T} \sum_{i \in \mathcal{T}}
  \left[\Delta Y_i - \sum_{j \in \mathcal{C}(i)} w_{ij} \cdot \Delta Y_j
  \right],
  \label{eq:psm-dd}
\end{equation}
où $\Delta Y_i = Y_{i,t=2} - Y_{i,t=1}$ est la variation intra-ménage
de la variable de résultat N-MODA, calculée directement sur les
ménages suivis entre les deux vagues. Le panel vrai dont nous disposons
rend cette différence directement observable sans agrégation par cellule,
ce qui renforce la validité de l'hypothèse de tendances parallèles. Les
écarts-types sont clusterisés au niveau de la grappe \parencite{caliendo2008}.

\subsubsection{Analyses d'hétérogénéité et de robustesse}
\label{subsec:robustesse}

Les analyses d'hétérogénéité estiment l'impact PSM-DD séparément selon le
milieu de résidence (urbain et rural), le groupe d'âge, le genre et chaque
dimension de privation. La robustesse est vérifiée par la comparaison entre les
trois algorithmes d'appariement, un test placebo et l'analyse de l'attrition.

La principale vérification de robustesse repose sur une approche par
panel d'enfants. La mesure N-MODA attribue à chaque enfant un jeu de
dimensions dépendant de son groupe d'âge (0-4, 5-14, 15-17 ans). Or un enfant
suivi entre 2018 et 2021 vieillit d'environ trois ans et peut donc changer de
groupe d'âge~: un enfant de 3 ans en 2018 (groupe 0-4) relève du groupe 5-14 en
2021, un enfant de 13 ans (groupe 5-14) passe au groupe 15-17. Les dimensions
et indicateurs qui lui sont applicables se modifient alors d'une vague à
l'autre. Plutôt que de comparer des enfants distincts d'une même tranche d'âge
à chaque vague, l'approche par panel d'enfants suit les mêmes enfants et
recalcule leur statut de privation N-MODA en tenant compte de cette transition
d'âge. Elle constitue une alternative au panel de ménages~: elle vérifie que
l'impact estimé des transferts sur la pauvreté multidimensionnelle n'est pas un
artefact de la composition par âge de l'échantillon, mais reflète bien une
évolution des privations des enfants réellement suivis.
